%=============================================================================
\section{Problem Setup}
\label{sec:setup}

Let $c$ denote the candidate system being evaluated. Let $(U, V)$ be the
pair of objects whose statistical dependence defines the reference-free signal.
In the simplest prompt--response form, $U = \text{prompt}$ and $V =
\text{response}$. In peer-prediction forms \citep{robertson2025limitations},
$(U, V)$ can be two reports tied to the same latent source.

\begin{definition}[Latent TVD-MI target]
\label{def:tvd-mi-target}
For candidate $c$, define
\begin{equation}
\tau(c)
:= \tvd\!\big(P^{(c)}_{UV},\; P^{(c)}_{U} \otimes P^{(c)}_{V}\big)
= \sup_{h \in \cH} \Delta_c(h),
\end{equation}
where $\cH = \{h : \cU \times \cV \to [0,1]\}$ and
\begin{equation}
\Delta_c(h)
:= \E_{P^{(c)}_{UV}}\!\big[h(U,V)\big]
   - \E_{P^{(c)}_{U} \otimes P^{(c)}_{V}}\!\big[h(U,V)\big].
\end{equation}
\end{definition}

\begin{definition}[Gold metric and validity link]
\label{def:gold-metric}
Let $g(c)$ be an inaccessible gold metric (e.g., human preference).
A \emph{monotone validity link} holds if
\begin{equation}
\tau(c_1) > \tau(c_2) \;\Longleftrightarrow\; g(c_1) > g(c_2).
\end{equation}
More generally, $g(c) = \phi(\tau(c))$ for a monotone function $\phi$.
\end{definition}

\begin{remark}
The architectural point is a two-layer decomposition: first prove consistency
of the ensemble estimator to~$\tau(c)$, then invoke the validity link to
transfer consistency to the gold metric~$g(c)$. This separation allows the
statistical analysis (Sections~\ref{sec:coverage}--\ref{sec:two-axis}) to
proceed independently of the domain-specific validity assumption
(Section~\ref{sec:gold}).
\end{remark}

\begin{remark}[Clarification on ``reference-free'']
Reference-free means \emph{no gold output}, not \emph{no source variable}.
The MI construction requires either prompt/source information or
same-vs-different pairing structure. Without structural assumptions,
elicitation without verification is impossible in full generality, so the
theory requires explicit identifiability assumptions rather than an
unconditional convergence claim.
\end{remark}
